\chapter{Produzione e scorte multiperiodali}
\label{cap:produzione}

\noindent\textbf{Classe:} LP / QP convesso \hfill
\textbf{Script:} \texttt{python/lab04\_produzione.py}

\section{Motivazione gestionale}

Un'azienda decide quanto produrre oggi e quanto conservare a scorta per la domanda
futura. Produrre in anticipo costa giacenza; produrre all'ultimo momento rischia di
sbattere contro il limite di capacità proprio nei mesi di picco. Il modello bilancia
margini, capacità, scorte e stabilità del piano --- ed è il punto di partenza ideale del
laboratorio perché ogni suo numero ha un'interpretazione economica immediata.



In pratica questo problema si presenta ogni volta che un'azienda pianifica la
produzione su più periodi con risorse condivise: piani principali di produzione
(MPS), pianificazione aggregata, contratti di fornitura con profili mensili. Il
nucleo della difficoltà è il \emph{legame intertemporale}: la decisione di oggi
(produrre in anticipo o no) cambia i vincoli di domani.

\textbf{In questo capitolo impareremo a:} (1) formulare un LP multiperiodale con
vincoli di bilancio delle scorte; (2) leggere e verificare i prezzi ombra della
capacità e i loro intervalli di validità; (3) usare un termine quadratico per
stabilizzare il piano; (4) costruire la curva del valore della capacità.

\section{Costruzione guidata del modello}

\paragraph{Il problema a parole.} \emph{Decidiamo} quante unità di ogni prodotto
produrre in ogni mese e quante tenerne a scorta. \emph{L'obiettivo} è minimizzare la
somma dei costi di produzione e di giacenza sull'intero orizzonte. \emph{I vincoli}
sono due: la contabilità di magazzino deve tornare in ogni mese (ciò che entra
uguaglia ciò che esce, e tutta la domanda va servita) e le ore macchina usate in un
mese non possono superare quelle disponibili.

\subsection*{Insiemi e dati (input del modello)}
\begin{center}\small
\begin{tabular}{llp{7.6cm}}
\toprule
Simbolo & Tipo & Significato \\
\midrule
$I$, $T$ & insiemi finiti & prodotti e periodi (mesi) \\
$d_{it}$ & $\in \R_{\ge 0}$ & domanda del prodotto $i \in I$ nel mese $t \in T$ \\
$c_{it},\, h_{it}$ & $\in \R_{\ge 0}$ & costo unitario di produzione e di giacenza \\
$a_i$ & $\in \R_{> 0}$ & ore macchina per unità del prodotto $i$ \\
$b_t$ & $\in \R_{\ge 0}$ & ore macchina disponibili nel mese $t$ \\
$s_{i0}$ & $\in \R_{\ge 0}$ & scorta iniziale del prodotto $i$ \\
\bottomrule
\end{tabular}
\end{center}

\subsection*{Variabili decisionali}
\begin{center}\small
\begin{tabular}{llp{7.6cm}}
\toprule
$x_{it}$ & $\ge 0$, continua & quantità prodotta di $i \in I$ nel mese $t \in T$ \\
$s_{it}$ & $\ge 0$, continua & scorta di $i$ alla fine del mese $t$ \\
\bottomrule
\end{tabular}
\end{center}

\begin{modello}[LP di produzione e scorte a costo minimo]
\begin{align}
\min\;& \sum_{i \in I}\sum_{t \in T} \bigl( c_{it}\, x_{it} + h_{it}\, s_{it} \bigr) \notag\\
\text{s.t.}\;\;
& s_{i,t-1} + x_{it} = d_{it} + s_{it} && \forall i \in I,\ t \in T \label{eq:bilancio}\\
& \sum_{i \in I} a_i\, x_{it} \le b_t && \forall t \in T \label{eq:capacita}\\
& x_{it},\, s_{it} \ge 0 && \forall i \in I,\ t \in T \notag
\end{align}
\end{modello}

\begin{spiegazione}
\textbf{Funzione obiettivo.} Somma, su tutti i prodotti e i mesi, i costi di
produzione $c_{it} x_{it}$ e i costi di giacenza $h_{it} s_{it}$: il modello cerca il
piano complessivamente più economico, libero di scambiare produzione anticipata
(più giacenza) con produzione al limite di capacità.

\textbf{Vincolo \eqref{eq:bilancio} (bilancio delle scorte).} È la ``contabilità di
magazzino'': ciò che entra (scorta precedente + produzione) è uguale a ciò che esce
(domanda servita + scorta successiva). È il vincolo che \emph{lega i mesi tra loro}: senza
di esso il problema si spezzerebbe in $|T|$ problemi indipendenti e le scorte non
avrebbero senso. Per $t=1$ si usa la scorta iniziale $s_{i0}$.

\textbf{Vincolo \eqref{eq:capacita} (capacità).} Le ore macchina sono condivise tra i
prodotti: è la risorsa scarsa che rende il problema interessante. Il suo prezzo ombra
dirà quanto vale un'ora in più.

\textbf{Cosa succede se\dots} togliamo il bilancio? Il modello produce zero (costo
minimo!) perché nulla lo obbliga a servire la domanda. Se la domanda totale supera la
capacità cumulata, il modello diventa inammissibile: in tal caso o si consente la
domanda persa con una variabile $u_{it} \ge 0$ e una penalità $\pi_{it}$ in obiettivo,
oppure si rivede la capacità.
\end{spiegazione}

\section{Esempio svolto a mano}

\begin{esempio}[Un prodotto, due mesi]
Domanda $d_1 = 60$, $d_2 = 100$; capacità $b_1 = b_2 = 80$ unità/mese ($a = 1$);
costo di produzione $c = 10$ \euro, giacenza $h = 1$ \euro/unità/mese, scorta iniziale nulla.

\textbf{Passo 1 — è ammissibile?} Domanda totale $160 \le$ capacità totale $160$: sì,
ma solo producendo al massimo in entrambi i mesi.

\textbf{Passo 2 — bilanci.} $x_1 = 60 + s_1$ e $s_1 + x_2 = 100$. Con $x_2 \le 80$
serve $s_1 \ge 20$, quindi $x_1 \ge 80$: ma $x_1 \le 80$, dunque $x_1 = 80$, $s_1 = 20$,
$x_2 = 80$.

\textbf{Passo 3 — costo.} $10\cdot(80+80) + 1\cdot 20 = 1620$~\euro. La soluzione è
\emph{forzata}: entrambi i vincoli di capacità sono attivi.

\textbf{Passo 4 — prezzo ombra a mano.} Con un'ora in più nel mese 2
($b_2 = 81$): $x_2 = 81$, $s_1 = 19$, $x_1 = 79$: il costo scende di $1$~\euro{}
(un'unità non sta più a scorta un mese). Il duale della capacità del mese 2 è
dunque $-1$~\euro/unità. Con un'ora in più nel mese 1 invece non cambia nulla
(il mese 1 produce già il necessario): duale $0$.
Questo meccanismo --- \emph{la capacità di un mese vale quanto giacenza fa risparmiare} ---
riappare identico nel caso di studio.
\end{esempio}

\section{Caso di studio}

Tre prodotti (indicati con $i = 1, 2, 3$) su sei mesi, con domanda stagionale crescente (picco nei mesi
4--5), un'unica risorsa condivisa e questi parametri (file CSV
\texttt{produzione\_domanda} e \texttt{produzione\_capacita} nella cartella
\texttt{dati/}):

\begin{center}\small
\begin{tabular}{lrrrr}
\toprule
Prodotto & $c_i$ (\euro/u) & $h_i$ (\euro/u/mese) & $a_i$ (ore/u) & $s_{i0}$ \\
\midrule
1 & 12{,}00 & 0{,}80 & 0{,}9 & 30 \\
2 & 18{,}00 & 1{,}20 & 1{,}4 & 20 \\
3 & 25{,}00 & 1{,}60 & 2{,}1 & 10 \\
\bottomrule
\end{tabular}
\qquad
\begin{tabular}{lrrrrrr}
\toprule
mese & 1 & 2 & 3 & 4 & 5 & 6 \\
\midrule
$b_t$ (ore) & 420 & 420 & 460 & 460 & 460 & 420 \\
$d_{1,t}$ & 110 & 130 & 150 & 190 & 210 & 160 \\
$d_{2,t}$ & 70 & 80 & 110 & 140 & 150 & 100 \\
$d_{3,t}$ & 50 & 55 & 60 & 80 & 95 & 70 \\
\bottomrule
\end{tabular}
\end{center}

\section{Implementazione}

Il cuore dello script (versione integrale in \texttt{lab04\_produzione.py}):

\begin{lstlisting}[style=python]
m = gp.Model("produzione_scorte")
x = m.addVars(prodotti, mesi, name="x")   # quantità prodotta
s = m.addVars(prodotti, mesi, name="s")   # scorta a fine mese

# bilancio: scorta entrante + produzione
#           = domanda + scorta uscente
m.addConstrs(
    ((s0[i] if t == 1 else s[i, t-1]) + x[i, t]
     == domanda[i, t] + s[i, t]
     for i in prodotti for t in mesi), name="bilancio")

# capacità condivisa (conservo i vincoli per i duali)
v_cap = m.addConstrs(
    (gp.quicksum(a[i] * x[i, t] for i in prodotti)
     <= capacita[t] for t in mesi), name="capacita")

m.setObjective(
    gp.quicksum(c[i]*x[i, t] + h[i]*s[i, t]
                for i in prodotti for t in mesi),
    GRB.MINIMIZE)
m.optimize()
\end{lstlisting}

\section{Risultati}

\begin{lstlisting}[style=output]
Costo totale ottimo: 32.889,33 EUR

Piano di produzione (unita'/mese):
mese       1      2      3      4      5      6
prod.1  80.0  130.0  150.0  190.0  210.0  160.0
prod.2  50.0   80.0  110.0  140.0  150.0  100.0
prod.3  89.5   91.0   81.4   44.3   29.0   64.8

Scorte a fine mese:
mese       1     2      3     4    5    6
prod.1   0.0   0.0    0.0   0.0  0.0  0.0
prod.2   0.0   0.0    0.0   0.0  0.0  0.0
prod.3  49.5  85.5  106.9  71.2  5.2  0.0
\end{lstlisting}

\begin{center}
\begin{tikzpicture}
\begin{axis}[lab, width=0.98\textwidth, height=6.2cm, ybar, bar width=5.5pt,
  title={Piano di produzione ottimo e domanda totale},
  xlabel={mese}, ylabel={unità}, ymin=0,
  xtick={1,2,3,4,5,6}, enlarge x limits=0.1, area legend]
\addplot [fill=teal!85, draw=none] table [col sep=comma, x=mese, y=x1] {\figdat{cap04_piano}};
\addplot [fill=rossomattone!80, draw=none] table [col sep=comma, x=mese, y=x2] {\figdat{cap04_piano}};
\addplot [fill=arancio!85, draw=none] table [col sep=comma, x=mese, y=x3] {\figdat{cap04_piano}};
\addplot [sharp plot, grigiomedio, thick, dashed, mark=o]
  table [col sep=comma, x=mese, y=domtot] {\figdat{cap04_piano}};
\legend{prodotto 1, prodotto 2, prodotto 3, domanda totale}
\end{axis}
\end{tikzpicture}
\captionof{figure}{Piano di produzione ottimo dell'LP: barre $=$ unità prodotte dei
tre prodotti in ciascun mese; linea tratteggiata $=$ domanda totale del mese. Nei mesi
di picco (4--5) la produzione totale resta sotto la domanda: la differenza è coperta
dalle scorte del prodotto 3 accumulate nei primi mesi.}
\end{center}

\begin{center}
\begin{tikzpicture}
\begin{axis}[labmezza, title={Scorte di fine mese}, xlabel={mese}, ylabel={unità}, ymin=0]
\addplot+ [mark=*] table [col sep=comma, x=mese, y=s1] {\figdat{cap04_piano}};
\addplot+ [mark=square*] table [col sep=comma, x=mese, y=s2] {\figdat{cap04_piano}};
\addplot+ [mark=triangle*] table [col sep=comma, x=mese, y=s3] {\figdat{cap04_piano}};
\legend{prodotto 1, prodotto 2, prodotto 3}
\end{axis}
\end{tikzpicture}%
\begin{tikzpicture}
\begin{axis}[labmezza, title={Effetto dello smoothing ($\gamma$)},
  xlabel={mese}, ylabel={produzione totale}]
\addplot+ [mark=*] table [col sep=comma, x=mese, y=gzero] {\figdat{cap04_smoothing}};
\addplot+ [mark=square*] table [col sep=comma, x=mese, y=gmezzo] {\figdat{cap04_smoothing}};
\addplot+ [mark=triangle*] table [col sep=comma, x=mese, y=gdue] {\figdat{cap04_smoothing}};
\legend{$\gamma=0$, $\gamma=0.5$, $\gamma=2$}
\end{axis}
\end{tikzpicture}
\captionof{figure}{A sinistra: scorte di fine mese per prodotto --- solo il prodotto 3 usa il
magazzino (strategia di pre-build), e le scorte si azzerano nel mese finale.
A destra: produzione totale mensile per tre valori del peso di smoothing $\gamma$:
già con $\gamma = 0{,}5$ il profilo diventa quasi piatto.}
\end{center}

\begin{insight}
Il solver ha scoperto da solo la strategia del \emph{pre-build}: i prodotti 1 e 2 (margini di
giacenza alti rispetto alle ore assorbite) vengono prodotti su domanda, mentre il
prodotto 3 --- il più oneroso in ore macchina --- viene prodotto in anticipo nei mesi
1--3 e accumulato (fino a 107 unità a scorta) per liberare capacità nei mesi di picco.
Le scorte tornano a zero nel mese finale: una scorta terminale positiva sarebbe denaro
sprecato, e il modello lo sa.
\end{insight}

\section{Analisi di sensitività}

\subsection*{Prezzi ombra della capacità}

\begin{lstlisting}[style=output]
mese 1: uso 330.0/420  | Pi =  0.000 EUR/ora | valido per b_t in [330.0, +inf]
mese 2: uso 420.0/420  | Pi = -0.762 EUR/ora | valido per b_t in [330.0, 524.0]
mese 3: uso 460.0/460  | Pi = -1.524 EUR/ora | valido per b_t in [370.0, 564.0]
mese 4: uso 460.0/460  | Pi = -2.286 EUR/ora | valido per b_t in [370.0, 564.0]
mese 5: uso 460.0/460  | Pi = -3.048 EUR/ora | valido per b_t in [399.0, 564.0]
mese 6: uso 420.0/420  | Pi = -3.810 EUR/ora | valido per b_t in [330.0, 431.0]

Verifica: +1 ora nel mese 6 -> costo da 32889.33 a 32885.52 (delta -3.810)
\end{lstlisting}

\begin{spiegazione}[Da dove viene il numero 0{,}762]
I duali crescono di esattamente $0{,}762$ \euro{} al mese. Non è un caso:
$0{,}762 = h_3 / a_3 = 1{,}6 / 2{,}1$. Un'ora in più nel mese $t$ permette di produrre
$1/a_3$ unità del prodotto 3 \emph{un mese più tardi}, risparmiando un mese di giacenza:
$h_3 \cdot (1/a_3)$. Il duale del mese 6 vale $5 \times 0{,}762 = 3{,}81$ perché
un'ora lì evita cinque mesi di giacenza a un'unità prodotta nel mese 1. Quando il
solver e un ragionamento economico elementare danno lo stesso numero, il modello è
probabilmente giusto.
\end{spiegazione}

\subsection*{Il valore della capacità e il costo della regolarità}

\begin{center}
\begin{tikzpicture}
\begin{axis}[labmezza, title={Curva valore della capacità},
  xlabel={capacità (\% dell'attuale)}, ylabel={costo ottimo (\euro)}]
\addplot+ [mark=*, mark size=1.4pt] table [col sep=comma, x=percento, y=costo]
  {\figdat{cap04_capacita}};
\addplot [arancio, dashed] coordinates {(100, 32300) (100, 33100)};
\end{axis}
\end{tikzpicture}%
\hfill
\begin{minipage}[b]{0.46\textwidth}\small
La curva è \textbf{convessa e decrescente}: le prime ore extra valgono molto, poi
sempre meno. Sotto il 97\% dell'attuale capacità il modello diventa
\emph{inammissibile}: la domanda non è più servibile per intero --- informazione
gestionale preziosa quanto l'ottimo stesso.

\vspace{1ex}
La variante QP aggiunge in obiettivo $\gamma \sum_{t \in T,\, t>1} (X_t - X_{t-1})^2$ con
$X_t = \sum_{i \in I} x_{it}$: con $\gamma = 0{,}5$ il profilo produttivo diventa quasi piatto
(variazione massima da 81 a 1 unità/mese) al costo di appena
$33.003 - 32.889 = 114$~\euro{} ($+0{,}35\%$).
\end{minipage}
\captionof{figure}{Costo totale ottimo in funzione della capacità disponibile
(in \% dell'attuale; la linea tratteggiata arancione segna la capacità corrente):
curva convessa e decrescente, interrotta a sinistra dove il problema diventa
inammissibile.}
\end{center}

\begin{insight}
Raccomandazione tipo da relazione di laboratorio: ``Conviene acquistare ore
straordinarie nei mesi 5--6 fino a 3--3,8 \euro/ora sotto il prezzo interno; un piano
livellato costa solo lo 0,35\% in più e riduce drasticamente i cambi di ritmo; sotto
il 97\% della capacità attuale il servizio completo non è garantibile.''
\end{insight}

\section{Esercizi}

\begin{esercizio}[verifica]
Verificare per perturbazione i duali dei mesi 3 e 4 (aggiungere un'ora e riottimizzare).
Perché il duale del mese 1 è nullo?
\end{esercizio}

\begin{esercizio}[domanda persa]
Aggiungere la variabile $u_{it}$ (domanda non servita) con penalità
$\pi = 40$ \euro/unità e ridurre la capacità al 90\%: quanta domanda si perde, di
quale prodotto, in quali mesi? Interpretare la scelta del modello.
\end{esercizio}

\begin{esercizio}[massimo profitto]
Trasformare il modello in massimizzazione del profitto con ricavi
$p = (20, 30, 42)$ \euro{} e servizio facoltativo ($y_{it} \le d_{it}$).
Quale domanda conviene \emph{non} servire, e perché?
\end{esercizio}

\begin{esercizio}[emissioni]
Ogni unità del prodotto 3 emette 2{,}5 kg di CO$_2$, i prodotti 1 e 2 un kg. Aggiungere il termine
$\tau \sum_{i \in I}\sum_{t \in T} e_i\, x_{it}$ all'obiettivo e tracciare la frontiera costo-emissioni per
$\tau \in [0, 10]$ \euro/kg (si veda l'analoga analisi nel
capitolo~\ref{cap:supplychain}).
\end{esercizio}

\begin{esercizio}[scorta di sicurezza]
Imporre $s_{it} \ge 0{,}1\, d_{i,t+1}$ (scorta di sicurezza del 10\%). Quanto costa
questa politica? Quale vincolo di capacità diventa più critico?
\end{esercizio}
